\lecture{6}
\title{Transformers}

\begin{document}

\subsection{CNN Limitations}

One limitation of CNNs is that they are not well-suited for
modeling \emph{long-distance} relationships.

\begin{center}
    \frame{\includegraphics[width=0.5\textwidth]{lecture-06/CNN-downside}}
    \emph{Question: How many birds are in this image?}
\end{center}

We'll introduce three architectural innovations to fix these problems.

\subsection{Tokens}

\textbf{Definition.}
A \textbf{token} is just a vector of neurons.
Connotatively, a token is an \emph{encapsulated} bundle of information.
You can tokenize \emph{anything}, whether it's image, text, or audio data.

Here are some details on how tokenization is implemented;
the general principle, though is that the details do not matter.
\begin{itemize}
    \item Text is tokenized by using a $\approx$ 100K-entry vocabulary.
    \begin{itemize}
        \item There's usually a pre-tokenization step in which
        a human-designed regex handles whitespace, punctuation, etc.
        \item The vocabulary is usually learned separately
        by simply reading text and identifying common bigrams, trigrams, etc.
    \end{itemize}
    \item Images are tokenized by simply partitioning the image
    into patches and vectorizing each patch.
    \begin{itemize}
        \item So a $224 \times 224$ image becomes $196$ vectors
        if patches of size $16 \times 16$ are used.
        \item Notably, the tokenization process does not take into account
        the ``content'' of the image.
    \end{itemize}
    \item Audio is tokenized by first compressing (e.g., via spectrogram),
    and only then patching and vectorizing.
\end{itemize}

\begin{center}
    \frame{\includegraphics[width=0.7\textwidth]{lecture-06/tokenization}}
\end{center}

Notably, a very simple tokenization method generally suffices,
because the NN on tokens, itself, can learn some of the
more interesting relationships between those tokens.

Here's how tokens look when used as nodes in a neural network (at left).
\begin{center}
    \frame{\includegraphics[width=0.6\textwidth]{lecture-06/token-NN}}
\end{center}
The linear portions of the NN work how you would expect them to.
\begin{center}
    \frame{\includegraphics[width=0.3\textwidth]{lecture-06/linear-token}}
\end{center}
The nonlinear mappings from individual tokens to individual tokens
are implemented not with a $\mathrm{ReLU}$, but with an NN! (Meta.)
\begin{center}
    \frame{\includegraphics[width=0.7\textwidth]{lecture-06/nonlinear-token}}
\end{center}

\subsection{Query-Key-Value Attention}

Currently, the architecture of NN has no way of identifying
which tokens are relevant and which tokens are not.
We would like our NN architecture to weight our tokens
based on how relevant they are to the query.

\begin{center}
    \frame{\includegraphics[width=0.6\textwidth]{lecture-06/attention-goal}}
\end{center}

One way to implement this is via \textbf{query-key-value attention}.
\begin{itemize}
    \item Each ``input token'' $\mathbf{t}$
    is mapped to a \textbf{query} token
    $\mathbf{q} := \mathbf{W}_q\mathbf{t}$
    (shown in orange, below).
    \item Each ``image token'' $\mathbf{t}$ 
    is mapped to a \textbf{key} token
    $\mathbf{k} := \mathbf{W}_k\mathbf{t}$
    (shown in purple, below).
    \item For each pair of image token $\mathbf{a}$
    and input token $\mathbf{b}$, we can determine
    the relevance of $\mathbf{a}$ and $\mathbf{b}$
    by computing the dot product of their query and key tokens
    $\mathbf{s} = \mathbf{q}_{\mathbf{b}}^{\top} \mathbf{k}_{\mathbf{a}}$.
    \item Each ``image token'' $\mathbf{t}$
    is mapped to a \textbf{value} token
    $\mathbf{v} := \mathbf{W}_v \mathbf{t}$.
    \begin{itemize}
        \item One could simply take $\mathbf{W}_v$ to be the identity;
        it's just extra expressivity.
    \end{itemize}
    \item The output $\mathbf{T}_{\mathrm{out}}$ is determined
    by a weighted sum of the values $\{\mathbf{v}\}$
    with weights given by the softmax-normalized scores $\mathbf{s}$.
\end{itemize}

\begin{center}
    \frame{\includegraphics[width=0.8\textwidth]{lecture-06/query-key-value}}
\end{center}

\subsection{Self-Attention}

One can also consider \textbf{self-attention}---instead of
computing the relevance between input and image,
we compute the relevance between some data and itself.

\begin{center}
    \frame{\includegraphics[width=0.7\textwidth]{lecture-06/self-attention}}
\end{center}

Notably, the map from $\mathbf{T}_{\mathrm{in}}$ to $\mathbf{T}_{\mathrm{out}}$
looks like:
\[
    \mathbf{T}_{\mathrm{out}}
    = \underbrace{\mathrm{softmax}\left(
        \frac{\mathbf{Q}_{\mathrm{in}}\mathbf{K}_{\mathrm{in}}^{\top}}{\sqrt{m}}
    \right)}_{\mathbf{A}} \mathbf{V}_{\mathrm{in}}
\]
where $m$ is the query/key dimension.
Note that $\mathbf{A}$ is normalized so that each of its rows sum to $1$,
so every output token is a weighted average of all the value tokens;
the weight on $\mathbf{v}_j$ in output token $i$ is
$\mathbf{A}_{ij} \propto \exp(\mathbf{q}_i^{\top}\mathbf{k}_j / \sqrt{m})$.
\begin{itemize}
    \item Relevance is measured under the learned maps:
    $\mathbf{q}_{\mathbf{a}}^{\top}\mathbf{k}_{\mathbf{b}}
    = \mathbf{a}^{\top}\mathbf{W}_q^{\top}\mathbf{W}_k\mathbf{b}$,
    which is \emph{not} symmetric.
    \begin{itemize}
        \item So ``A attends to B'' is distinct from
        ``B attends to A'' (and neither are exactly equivalent to
        ``A and B are similar'').
    \end{itemize}
    \item Unlike in CNNs,
    every pair of tokens is connected by a single matrix multiplication,
    no matter how far apart they are.
\end{itemize}
Below, the learned projections happen to make each bird-head patch
attend to the other bird-head patches.
\begin{center}
    \frame{\includegraphics[width=0.6\textwidth]{lecture-06/turkey-transformer}}
\end{center}

\textbf{Definition.} There's also \textbf{multihead self-attention (MSA)},
in which we concatenate $k$ different attempts at self-attention tokenwise,
then project back down with a matrix $\mathbf{W}_{\mathtt{MSA}}$.

\begin{center}
    \frame{\includegraphics[width=0.4\textwidth]{lecture-06/MSA}}
\end{center}

\subsection{Transformer Architecture}

Putting together multihead self-attention (MSA) with
multi-layer perceptrons (MLPs), we get a transformer!

\begin{center}
    \frame{\includegraphics[width=0.3\textwidth]{lecture-06/transformer}}
\end{center}

The code for a transformer is extremely simple, too.

\begin{verbatim}[python]
# x: input data (RGB image)
# K: tokenization patch size
# d: token/query/key/value dimensionality (setting these all as the same)
# L : number of layers
# W_q_T, W_k_T, W_v_T : transposed query/key/value projection matrices
# mlp: tokenwise mlps

# tokenize input image
T = tokenize(x, K) # 3 x Hx W image --> N xd array of token code vectors

# run tokens through all L layers
for l in range(L):

    # attention layer
    Q, K, V = nn.matmul(nn.layernorm(T), [W_q_T[l], W_k_T[l], W_v_T[l]])
    # nn.matmul does matrix multiplication
    A = nn.softmax(nn.matmul(Q,K.transpose())/sqrt(d), dim=-1)
    T = nn.matmul(A, V) + T # note residual connection

    # tokenwise mlp
    T = mlp[l](nn.layernorm(T)) + T # note residual connection

# T now contains the output token representation computed by the transformer
\end{verbatim}

\subsection{Positional Encodings}

One observation about the transformer architecture
is that it is \textbf{permutation equivariant}.
\[ F_{\theta}(\pi(\mathbf{T}_{\mathrm{in}}))
= \pi(F_{\theta}(\mathbf{T}_{\mathrm{in}}))
~~~ \text{ and } ~~~
\textsc{attn}(\pi(\mathbf{T}_{\mathrm{in}}))
= \pi(\textsc{attn}(\mathbf{T}_{\mathrm{in}}))
~~~ \text{ for any MLP } F_{\theta}
~ \text{ and any MSA } \textsc{attn}. \]
But we actually \emph{don't} want permutation equivariance,
because that implies that it may be difficult for models to
learn about forms of data that care about position.
The fix, naturally, is to add position to each of our tokens.
\begin{center}
    \frame{\includegraphics[width=0.25\textwidth]{lecture-06/positional-encoding}}
\end{center}
There is some nuance concerning how, exactly, we encode position.
For images, for example, we might encode position
by the Fourier coefficients associated with that position.
\begin{center}
    \frame{\includegraphics[width=0.5\textwidth]{lecture-06/Fourier-encoding}}
\end{center}
Similarly, positional encodings of points on a sphere
might use spherical harmonics.
\begin{center}
    \frame{\includegraphics[width=0.2\textwidth]{lecture-06/spherical-harmonics}}
\end{center}
Positional encodings of nodes in a graph are given by the eigenvectors
of its Laplacian. (Details omitted.)
\begin{center}
    \frame{\includegraphics[width=0.5\textwidth]{lecture-06/graph-encoding}}
\end{center}

\subsection{Autoregressive Models (GPT, Claude)}

Recall that LLMs like GPT generate text by ``predicting the next word''.
\begin{center}
    \frame{\includegraphics[width=0.5\textwidth]{lecture-06/text-prediction}}
\end{center}
GPT's means of prediction is basically just a transformer.
The only distinction is that the attention matrix must forbid
some tokens attending to tokens that occur in the future.
To implement this, it suffices to mask $\mathbf{A}$ like so:
\begin{center}
    \frame{\includegraphics[width=0.2\textwidth]{lecture-06/masked-attention}}
\end{center}
More generally, here's how the causal self-attention relationships look
for autoregressive image-to-text architectures.
\begin{center}
    \frame{\includegraphics[width=0.5\textwidth]{lecture-06/image-to-text}}
\end{center}
We call this \emph{cross-modality attention}, because
attention occurs between data in both text and image modalities.

\end{document}